% 07_writing_effectively.tex — The Craft of Writing Effectively by Larry McEnerney

\section{The Craft of Writing Effectively by Larry McEnerney}
\label{sec:writing_effectively}

Larry McEnerney served as the director of the University of Chicago's academic writing program (colloquially known as the Little Red Schoolhouse) for several decades.
In a famous lecture, McEnerney discusses how graduate students are ill-prepared to write in academic settings despite having ``learned'' to write for years, because many of the rules that they have been taught are counterproductive for academia.
\textbf{He argues that writing is about \textit{readers} instead of writers.}

\textbf{Thinking and writing are not separate processes} like many students learn in grade school.
People with expert knowledge use writing to help themselves think, because their work is so complex that writing is necessary.
However, writers generate a text ``horizontally'' (as they think about the world) while readers interpret a text ``vertically'' (as they read about the world).
The language patterns used for writing \& thinking interfere with reading \& understanding, causing readers to slow down, become aggravated, and eventually stop.
In the past, students' readers were mostly their teachers, who continued reading because they were paid to do so; in academia, people are not paid to care.

\textbf{Good writing needs to provide value over anything else.}
Many people think that effective writing is characterized by its clarity, organization, or persuasiveness.
These are certainly important qualities---but as McEnerney declares, ``If it's clear and useless, it's useless.
If it's organized and useless, it's useless.''
Value comes from a community of readers, not inherently from content.

\textbf{Writing is not about communicating \textit{your} ideas: instead, it is about changing your \textit{readers'} ideas.}
Students are asked to explain their thoughts to demonstrate their understanding.
In academia, the question shifts from ``Why do \textit{you} think that?'' to ``Why should \textit{I} think that?''
Academic knowledge is not accepted until it has been properly challenged.

To that end, \textbf{just because something is original or new does not mean it is valuable.}
In a positivistic world, ``knowledge'' accumulates over time, and any new discoveries are added to the total.
In practice, ``knowledge'' actually consists of conversations moving through time, while a group of people determine what is included.
McEnerney acknowledges the inherent flaw that certain groups---which, historically, did not include many diverse voices---are given the power to make such determinations.
However, he concedes that it is how academia functions, and that scholars must understand the system, however flawed it may be, in order to advance their careers.
Two caveats to this model are (1) the group that determines ``knowledge'' is changing and expanding, and (2) the ``boundary'' of knowledge is permeable---wrong ideas are left behind.

\textbf{Every community has its own ``code'' or set of words that create value.}
Persuasion depends on a thorough understanding of what readers doubt in order to overcome those doubts.
To create value, writers need to challenge existing ideas within the community's code.
